\section{Equilibrium states and prime occupation spaces}
\label{sec:equilibrium}

The local integral gives a geometric series because the valuation takes
nonnegative integer values on the unit ball. A Hilbert space with one
basis vector for each such value gives the same series as a heat trace.
This second construction also has an observable algebra and a time
evolution. We first determine its equilibrium state, then add the phase
observables that distinguish the full arithmetic system.

All Hilbert spaces in this section are complex, with inner products linear
in the first variable. Their operator algebras use composition as product,
the Hilbert adjoint as involution, and the operator-norm topology.
They have cohomological degree zero; their frequency decomposition under
time evolution is an additional grading, not a cochain differential.

\subsection{The equilibrium condition}

Let \(A\) be a unital complex \(C^*\)-algebra and
\(\tau:\mathbb R\to\operatorname{Aut}(A)\) a strongly continuous group.
A state is a positive complex-linear functional \(\varphi:A\to\mathbb C\)
with \(\varphi(1)=1\). An eigenvector of frequency \(\omega\in\mathbb R\)
satisfies \(\tau_t(a)=e^{it\omega}a\). An element is entire analytic if
\(t\mapsto\tau_t(a)\) extends to an entire norm-holomorphic function.

\begin{definition}[KMS condition]\label{bc:def:kms}
For \(\beta>0\), a state \(\varphi\) is KMS\(_\beta\) if, for all
\(a,b\in A\), there is a bounded continuous function on the strip
\(0\leq\operatorname{Im}z\leq\beta\), holomorphic in its interior, such that
\[
 F_{a,b}(t)=\varphi(a\tau_t(b)),\qquad
 F_{a,b}(t+i\beta)=\varphi(\tau_t(b)a).
\]
For entire analytic elements the corresponding algebraic identity is
\begin{equation}\label{bc:eq:kms}
 \varphi(ab)=\varphi(b\tau_{i\beta}(a)).
\end{equation}
\end{definition}

\begin{lemma}[The analytic criterion]\label{bc:lem:analytic}
Suppose \(\mathcal B_\tau\) is a norm-dense star subalgebra of \(A\)
spanned by eigenvectors of \(\tau\). A state is KMS\(_\beta\) exactly when
\eqref{bc:eq:kms} holds for \(a,b\in\mathcal B_\tau\).
\end{lemma}

\begin{proof}
For \(a,b\in\mathcal B_\tau\), the function
\(\varphi(a\tau_z(b))\) is a finite sum of exponentials and is bounded
on the closed strip. The difference between this function and a KMS
strip function has zero lower boundary. Schwarz reflection across that
boundary, followed by the identity theorem, makes the difference zero.
Its upper boundary then gives \eqref{bc:eq:kms} after interchanging
\(a\) and \(b\).

Conversely, the identity gives both required boundaries for the finite
exponential sum. Their absolute values are at most \(\|a\|\|b\|\).
The bounded-strip maximum principle gives the same bound in the strip.
Approximate arbitrary \(a,b\) in norm by elements of \(\mathcal B_\tau\).
The same estimate applied to differences makes the strip functions
uniformly Cauchy. Their limit is continuous on the closed strip,
holomorphic in its interior, and has the required boundaries.
\end{proof}

\subsection{A prime oscillator}

Fix a prime \(p\).  Let \(\mathscr H_p=\ell^2(\mathbb N_0)\) with
orthonormal basis \((e_{p,r})_{r\geq0}\), and set
\[
 S_pe_{p,r}=e_{p,r+1},\qquad \mathcal T_p=C^*(S_p)
 \subset\mathcal B(\mathscr H_p).
\]
Thus \(S_p^*S_p=1\) but \(S_pS_p^*\ne1\).  Define the positive operator
\[
 N_pe_{p,r}=r\log p\,e_{p,r}
\]
on its maximal domain
\[
 \mathcal D(N_p)=
 \left\{\sum_{r\geq0}a_re_{p,r}:
 \sum_{r\geq0}r^2(\log p)^2|a_r|^2<\infty\right\}.
\]
Testing its adjoint against each basis vector gives exactly the same
maximal domain, so the diagonal operator is self-adjoint. Truncation of
the coordinate sequence converges in its graph norm; hence finite-support
vectors form a core.

\begin{proposition}[Prime heat trace and Toeplitz KMS state]
\label{prop:prime-kms}
The unitaries \(U_{p,t}=e^{itN_p}\) define a strongly continuous dynamics
\[
 \alpha_{p,t}(a)=U_{p,t}aU_{p,t}^*,\qquad
 \alpha_{p,t}(S_p)=p^{it}S_p
\]
on \(\mathcal T_p\).  For \(s\in\mathbb C\),
\(e^{-sN_p}\) is trace class exactly when \(\operatorname{Re}s>0\), and
then
\[
 \operatorname{Tr}_{\mathscr H_p}(e^{-sN_p})=(1-p^{-s})^{-1}.
\]
For every real \(\beta>0\), \((\mathcal T_p,\alpha_p)\) has the unique
KMS\(_\beta\) state
\[
 \varphi_{\beta,p}(a)
 =(1-p^{-\beta})
 \operatorname{Tr}_{\mathscr H_p}(e^{-\beta N_p}a).
\]
On the dense analytic star algebra
\(\mathcal P_p=\operatorname{span}\{S_p^mS_p^{*n}:m,n\geq0\}\),
\[
 \varphi_{\beta,p}(S_p^mS_p^{*n})
 =\delta_{m,n}p^{-\beta m}.
\]
\end{proposition}

\begin{proof}
The formula for the dynamics follows on the basis vectors and extends by
norm density from \(\mathcal P_p\).  The singular values of
\(e^{-sN_p}\) are \(p^{-r\operatorname{Re}s}\), so their sum is finite
exactly when \(\operatorname{Re}s>0\); summing the diagonal entries gives
the trace formula.

For \(q=p^{-\beta}\), the operator
\((1-q)e^{-\beta N_p}\) is positive and has trace one.  Every monomial in
\(\mathcal P_p\) is entire analytic, and
\[
 e^{-\beta N_p}a=\alpha_{p,i\beta}(a)e^{-\beta N_p}.
\]
Cyclicity of the trace gives
\(\varphi_{\beta,p}(ab)=
\varphi_{\beta,p}(b\alpha_{p,i\beta}(a))\), the KMS identity in this
convention.  Direct evaluation gives the monomial formula.

Conversely, applying \eqref{bc:eq:kms} with \(b=1\) makes every
nonzero-frequency monomial have expectation zero.  For the remaining monomials,
the KMS identity gives
\[
 \varphi(S_p^mS_p^{*m})
 =p^{-\beta m}\varphi(S_p^{*m}S_p^m)=p^{-\beta m}.
\]
Lemma~\ref{bc:lem:analytic} proves the strip condition, and density of
\(\mathcal P_p\) proves uniqueness.  These formulas agree with
the prime system of Bost and Connes
\cite[\S2, Propositions~7--8]{BostConnes1995}.
\end{proof}

\begin{remark}[Heat flow is not observable dynamics]
For \(t>0\), the assignment \(S_p\mapsto p^{-t}S_p\) cannot be a unital
star-homomorphism because
\[
 (p^{-t}S_p)^*(p^{-t}S_p)=p^{-2t}1\ne1.
\]
The observable evolution is the automorphism group
\(S_p\mapsto p^{it}S_p\).  The unnormalized partition trace is not a
state, since its value at the identity is \((1-p^{-\beta})^{-1}\).
\end{remark}

\subsection{The tensor product over the primes}

Write \(\mathbb N^\times=\{1,2,\ldots\}\). For a finite set \(F\) of
primes, put
\[
 \mathcal T_F=\bigotimes_{p\in F}^{\min}\mathcal T_p,
 \qquad
 \mathcal T_{\mathrm{fin}}=\varinjlim_F\mathcal T_F,
\]
where the transition maps are \(a\mapsto a\otimes1\), and use the product
dynamics.  Relative to the vacuum vectors \(e_{p,0}\), unique
factorization gives
\[
 \bigotimes_p^{e_{p,0}}\mathscr H_p\simeq\ell^2(\mathbb N^\times).
\]
The compatible tensor-product representations define
\[
 \pi_{\mathrm{fin}}:\mathcal T_{\mathrm{fin}}
 \longrightarrow\mathcal B(\ell^2(\mathbb N^\times)),
 \qquad \pi_{\mathrm{fin}}(S_p)e_n=e_{pn}.
\]
On the latter space define
\[
 He_n=(\log n)e_n,\qquad
 \mathcal D(H)=
 \left\{\sum_{n\geq1}a_ne_n:
 \sum_{n\geq1}(\log n)^2|a_n|^2<\infty\right\}.
\]

\begin{proposition}[The noninteracting prime subsystem]
\label{prop:prime-product}
For every \(\beta>0\), \((\mathcal T_{\mathrm{fin}},\alpha)\) has one
KMS\(_\beta\) state \(\Phi_\beta\), characterized by
\[
 \Phi_\beta|_{\mathcal T_F}
 =\bigotimes_{p\in F}\varphi_{\beta,p}.
\]
The operator \(H\) is positive and self-adjoint on the displayed domain,
with finite-support vectors as a core, and implements the dynamics in the precise sense
\[
 \pi_{\mathrm{fin}}(\alpha_t(a))
 =e^{itH}\pi_{\mathrm{fin}}(a)e^{-itH}
 \qquad(a\in\mathcal T_{\mathrm{fin}}).
\]
It satisfies
\[
 e^{-sH}\text{ is trace class}
 \quad\Longleftrightarrow\quad \operatorname{Re}s>1.
\]
In that half-plane,
\begin{equation}
 \operatorname{Tr}(e^{-sH})=\sum_{n\geq1}n^{-s}=\zeta(s).
\label{bc:eq:trace}
\end{equation}
For \(\beta>1\), but not for \(0<\beta\leq1\), the product state is the
normal Gibbs state
\[
 \Phi_\beta(a)=\zeta(\beta)^{-1}
 \operatorname{Tr}(\pi_{\mathrm{fin}}(a)e^{-\beta H}).
\]
\end{proposition}

\begin{proof}
On a finite tensor product, \eqref{bc:eq:kms} with \(b=1\) kills every
monomial except those with zero total frequency.  Unique factorization of integers makes zero
total frequency equivalent to zero frequency at every prime.  Repeated use
of the KMS identity then forces the product of the local values in
Proposition~\ref{prop:prime-kms}. Conversely, the tensor product of the
local states satisfies the KMS identity on products of local analytic
monomials, by applying the local identity in each factor.
Lemma~\ref{bc:lem:analytic} proves existence at each finite stage.
The states are compatible under \(a\mapsto a\otimes1\); they therefore
define a state on the norm-completed inductive limit, where the same
analytic criterion applies. This proves existence and uniqueness there.

Under the displayed unitary, the sum of the local energies is \(\log n\).
Testing the adjoint on the basis identifies its domain with
\(\mathcal D(H)\), and truncating a sequence converges in the graph norm.
The singular values of \(e^{-sH}\) are \(n^{-\operatorname{Re}s}\).  Their
sum is finite exactly for \(\operatorname{Re}s>1\).  Normalizing by this
sum gives the Gibbs formula in that half-plane.  Below it, the abstract
product KMS state still exists, but no trace-class density
\(e^{-\beta H}\) exists in this representation.
\end{proof}

\subsection{The universal phase algebra}

The prime shifts remember divisibility but do not distinguish the phases
of roots of unity. To retain these phases, adjoin the continuous characters
of the profinite integers. The resulting algebra has the same occupation
Hamiltonian, but a larger space of equilibrium states.

Let \(\bcN=\{1,2,\ldots\}\), and let \(\bcZhat=\varprojlim_N\mathbb Z/N\mathbb Z\)
with its compact inverse-limit topology. The Chinese remainder theorem
identifies this ring with the product \(\prod_p\mathbb Z_p\) used in the
adelic construction: compatible residues modulo every positive integer
are equivalent to compatible residues modulo every prime power.
Its unit group is
\(\bcUnits=\varprojlim_N(\mathbb Z/N\mathbb Z)^\times\).
For \(r=a/N\in\mathbb Q/\mathbb Z\) and \(x\in\bcZhat\), set
\[
 \bcphase_r(x)=\exp(2\pi i a x_N/N),
\]
where \(x_N\) is the residue of \(x\) modulo \(N\).
This definition is independent of the chosen presentation of \(r\).
The functions \(\bcphase_r\) separate points. Their complex linear span is
therefore uniformly dense in \(C(\bcZhat)\).

The universal unital complex \(C^*\)-algebra \(\bcA\) has isometries
\(\mu_n\), \(n\in\bcN\), and unitaries \(e(r)\),
\(r\in\mathbb Q/\mathbb Z\), with relations
\begin{equation}\label{bc:eq:relations}
\begin{gathered}
 \mu_n^*\mu_n=1,\qquad \mu_n\mu_m=\mu_{nm},\qquad
 \mu_n\mu_m^*=\mu_m^*\mu_n\quad(\gcd(n,m)=1),\\
 e(0)=1,\qquad e(r)^*=e(-r),\qquad e(r)e(s)=e(r+s),\\
 e(r)\mu_n=\mu_ne(nr),\qquad
 \mu_ne(r)\mu_n^*=\frac1n\sum_{n\delta=r}e(\delta).
\end{gathered}
\end{equation}
The sum in the last relation has exactly \(n\) terms. Concretely, take
the free unital star algebra on these symbols, impose the displayed
relations, and use the supremum of the operator norms over all star
representations by isometries and unitaries. Each polynomial has finite
supremum norm, bounded by the sum of the absolute values of its
coefficients. Quotient by the zero seminorm ideal and complete.
The representations below prove that this completion is nonzero.
The dynamics is
\begin{equation}\label{bc:eq:dynamics}
 \sigma_t(\mu_n)=n^{it}\mu_n,
 \qquad \sigma_t(e(r))=e(r),\qquad t\in\mathbb R.
\end{equation}
These are the presentation and time convention of
Bost--Connes \cite{BostConnes1995}*{Proposition~18 and \S5}.
The relations also imply \(\mu_1=1\): multiply
\(\mu_1^2=\mu_1\) on the left by \(\mu_1^*\).

\begin{proposition}[The phase representation]\label{bc:prop:representation}
For every \(u\in\bcUnits\), the formulas
\begin{equation}\label{bc:eq:representation}
 \pi_u(\mu_m)\varepsilon_k=\varepsilon_{mk},\qquad
 \pi_u(e(r))\varepsilon_k=\bcphase_r(ku)\varepsilon_k
\end{equation}
define a unital star representation of \(\bcA\) on
\(\mathscr H=\ell^2(\bcN)\). Its inner product is linear in the first variable.
The phase subalgebra
\(D=C^*(e(r):r\in\mathbb Q/\mathbb Z)\) is canonically \(C(\bcZhat)\),
with \(e(r)=\bcphase_r\).

For the operator \(H\) and its closed diagonal domain in
Proposition~\ref{prop:prime-product},
\[
 \pi_u(\sigma_t(a))=e^{itH}\pi_u(a)e^{-itH}.
\]
Its occupation trace is \eqref{bc:eq:trace}, with the exact domain
\(\operatorname{Re}s>1\).
\end{proposition}

\begin{proof}
The adjoint of the shift by \(m\) sends \(\varepsilon_k\) to
\(\varepsilon_{k/m}\) if \(m\mid k\), and to zero otherwise.
This proves the isometry and semigroup relations. If \(\gcd(m,n)=1\),
then \(m\mid nk\) if and only if \(m\mid k\). This proves the mixed
commutation relation. The phase multiplication and its covariance with
the shifts follow by evaluation on \(\varepsilon_k\).

For the last relation in \eqref{bc:eq:relations}, choose \(\delta_0\) with
\(n\delta_0=r\). Character orthogonality gives
\[
 \frac1n\sum_{n\delta=r}\bcphase_\delta(ku)
 =\bcphase_{\delta_0}(ku)\frac1n
      \sum_{j=0}^{n-1}\exp(2\pi i jku/n)
 =\begin{cases}
   \bcphase_r((k/n)u),&n\mid k,\\
   0,&n\nmid k.
  \end{cases}
\]
Here the finite sum uses the residue of \(u\) modulo \(n\).
Its invertibility makes its vanishing condition precisely \(n\mid k\).
This verifies all relations and also proves that the universal algebra
is nonzero. Its universal norm exists because every generator has norm
one in every representation.

The dual of the discrete group \(\mathbb Q/\mathbb Z\) is \(\bcZhat\).
Indeed, a character is specified by compatible values on \(1/N\),
which are precisely the residues \(x_N\). Thus \(D\) is a quotient
of \(C(\bcZhat)\). In \(\pi_1\), an element \(f\in C(\bcZhat)\) has
diagonal values \(f(k)\). The positive integers are dense in \(\bcZhat\),
so this representation is faithful on \(C(\bcZhat)\). The quotient has
zero kernel.

The implementation of \(\sigma\) holds on generators by evaluation on
basis vectors and then on the algebra by norm continuity. The diagonal
domain, core, and trace statement are those already proved in
Proposition~\ref{prop:prime-product}.
\end{proof}

At a finite set \(F\) of primes, the vectors indexed by
\(n=\prod_{p\in F}p^{k_p}\) have energies \(\sum_{p\in F}k_p\log p\).
Hence
\begin{equation}\label{bc:eq:finite-trace}
 \sum_{n\in\bcN_F}n^{-\beta}
   =\prod_{p\in F}(1-p^{-\beta})^{-1},\qquad \beta>0,
\end{equation}
where \(\bcN_F\) consists of positive integers with prime factors in \(F\).
For \(F=\{2,3\}\) and \(\beta=1\), this trace is \(3\).
The full trace diverges at the same inverse temperature. Classifying
equilibrium states therefore requires a construction beyond that trace.

\subsection{From the universal algebra to scaling measures}

For \(f\in C(\bcZhat)\) and \(n\in\bcN\), define
\begin{equation}\label{bc:eq:alpha}
 L_nf(x)=f(nx),\qquad
 \alpha_n(f)(x)=
 \begin{cases}f(x/n),&x\in n\bcZhat,\\0,&x\notin n\bcZhat.
 \end{cases}
\end{equation}
Multiplication by \(n\) is a homeomorphism from \(\bcZhat\) onto the
clopen subset \(n\bcZhat\). In particular, \(\alpha_n(f)\) is continuous.
The relations imply
\begin{equation}\label{bc:eq:covariance}
 f\mu_n=\mu_nL_nf,\qquad
 \mu_n^*f=(L_nf)\mu_n^*,\qquad
 \mu_nf\mu_n^*=\alpha_n(f).
\end{equation}
The last identity follows first for characters from the finite sum in
\eqref{bc:eq:relations}, and then for every continuous function by density.

\begin{lemma}[Normal forms and expectation]\label{bc:lem:expectation}
The linear span of \(\mu_mf\mu_n^*\), with \(m,n\in\bcN\) and
\(f\in D\), is a dense star subalgebra of \(\bcA\). It admits a
positive unital conditional expectation \(E:\bcA\longrightarrow D\)
with
\begin{equation}\label{bc:eq:expectation}
 E(\mu_mf\mu_n^*)=\begin{cases}\alpha_n(f),&m=n,\\0,&m\ne n.
 \end{cases}
\end{equation}
Every monomial can be written with coprime \(m,n\), after replacing
its coefficient in \(D\).
\end{lemma}

\begin{proof}
If \(d=\gcd(n,k)\), the isometry and coprime relations give
\(\mu_n^*\mu_k=\mu_{k/d}\mu_{n/d}^*\). Thus
\[
 (\mu_mf\mu_n^*)(\mu_kg\mu_l^*)
 =\mu_{mk/d}\big((L_{k/d}f)(L_{n/d}g)\big)\mu_{ln/d}^*.
\]
This proves closure under products. Adjoints and all generators also
belong to the span. If \(d=\gcd(m,n)\), then
\[
 \mu_mf\mu_n^*=\mu_{m/d}\alpha_d(f)\mu_{n/d}^*,
\]
which proves reduction to coprime indices.

The compact group \(K=\prod_{p\ \mathrm{prime}}\mathbb T\) acts by
\[
 \kappa_z(\mu_n)=z(n)\mu_n,\qquad
 z(n)=\prod_pz_p^{v_p(n)},\qquad \kappa_z(f)=f.
\]
All relations are preserved. The inverse is \(\kappa_{z^{-1}}\).
Norm continuity follows first on monomials and then by density.
Integrating this action against probability Haar measure on \(K\)
defines a positive unital map \(E\). Orthogonality of its characters
gives \eqref{bc:eq:expectation}. Its range is exactly \(D\), and the
identity \(E(fag)=fE(a)g\), for \(f,g\in D\), gives the conditional
expectation property.
\end{proof}

\begin{proposition}[Faithfulness and the prime subalgebra]
\label{bc:prop:faithful}
Each representation \(\pi_u\), \(u\in\bcUnits\), is faithful.
Its restriction identifies \(C^*(\mu_p:p\text{ prime})\) with
\(\mathcal T_{\mathrm{fin}}\), sending \(\mu_p\) to the prime shift.
\end{proposition}

\begin{proof}
The expectation \(E\) is faithful. Indeed, if \(a\geq0\) and \(E(a)=0\),
then for every state \(\omega\) the nonnegative continuous function
\(z\mapsto\omega(\kappa_z(a))\) has Haar integral zero. Haar measure has
full support, so its value at the identity is zero. States separate
positive elements, and hence \(a=0\).

On \(\ell^2(\bcN)\) let \(V_z\varepsilon_k=z(k)\varepsilon_k\).
The formula on generators gives
\(\pi_u(\kappa_z(a))=V_z\pi_u(a)V_z^*\).
If \(\pi_u(a)=0\), integration implies \(\pi_u(E(a^*a))=0\).
The set \(\{ku:k\geq1\}\) is dense in \(\bcZhat\), because the
positive integers are dense and multiplication by \(u\) is a
homeomorphism. Thus \(\pi_u\) is faithful on \(D\), so
\(E(a^*a)=0\), and faithfulness of \(E\) gives \(a=0\).

The shifts \(\pi_u(\mu_p)\) do not depend on \(u\).
For a finite set \(F\), unique factorization decomposes
\(\ell^2(\bcN)\) into the tensor product of the occupation spaces at
\(p\in F\) and the space indexed by integers with no prime factor in
\(F\). The shifts act on the first factor and as the identity on the
second. This is precisely the faithful spatial representation of
\(\bigotimes_{p\in F}^{\min}\mathcal T_p\).
Taking the norm closure of the increasing finite prime subalgebras
proves the claimed identification.
\end{proof}

The same normal forms with coefficients in \(\operatorname{span}\{e(r)\}\)
give a dense entire analytic star algebra \(\mathcal B\). Its monomials
have frequency \(\log(m/n)\). Preservation of the defining relations
also proves that \eqref{bc:eq:dynamics} extends to automorphisms, with
inverse \(\sigma_{-t}\). The dynamics is strongly continuous by density.

Lemma~\ref{bc:lem:analytic} applies to this algebra with \(\tau=\sigma\).

\begin{theorem}[The measure correspondence]\label{bc:thm:measures}
Fix \(\beta>0\). Restriction to \(D=C(\bcZhat)\) gives an affine bijection
between the KMS\(_\beta\) states on \(\bcA\) and the probability
Borel measures \(\nu\) on \(\bcZhat\) satisfying
\begin{equation}\label{bc:eq:scaling}
 \nu(nB)=n^{-\beta}\nu(B)
 \quad(n\in\bcN,\ B\subset\bcZhat\text{ Borel}).
\end{equation}
The inverse sends \(\nu\) to
\(\varphi_\nu(a)=\int_{\bcZhat}E(a)\,d\nu\).
\end{theorem}

\begin{proof}
For a monomial \(a\) of frequency \(\omega\), the KMS identity with
\(b=1\) gives \(\varphi(a)=e^{-\beta\omega}\varphi(a)\).
Thus \(\varphi\) vanishes on every nonzero frequency and
\(\varphi=\varphi|_D\circ E\). Riesz representation identifies
\(\varphi|_D\) with a probability measure \(\nu\).
The KMS identity applied to \(a=\mu_n\) and \(b=f\mu_n^*\) gives
\[
 \int\alpha_n(f)\,d\nu=n^{-\beta}\int f\,d\nu.
\]
This holds for continuous \(f\) by norm density. Uniqueness of regular
Borel measures and the homeomorphism \(x\mapsto nx\) identify it
with \eqref{bc:eq:scaling}.

Conversely, \(\varphi_\nu\) is a state because \(E\) is positive and
unital. It remains to verify \eqref{bc:eq:kms} on monomials.
Products of two monomials have zero expectation unless the sum of
their frequencies is zero. In that case, reduction to
coprime indices writes the two monomials as
\(a=\mu_mf\mu_n^*\) and \(b=\mu_ng\mu_m^*\). Consequently,
\[
 \varphi_\nu(ab)=m^{-\beta}\int fg\,d\nu,
 \qquad
 \varphi_\nu(ba)=n^{-\beta}\int gf\,d\nu.
\]
Since \(\sigma_{i\beta}(a)=(m/n)^{-\beta}a\), these equations are
exactly \eqref{bc:eq:kms}. In the other cases both sides vanish.
Lemma~\ref{bc:lem:analytic} completes the proof.
\end{proof}

\subsection{Finite-prime partitions}

For a finite set \(F\) of primes, put
\[
 W_F=\bcZhat\setminus\bigcup_{p\in F}p\bcZhat,
 \qquad c_F(\beta)=\prod_{p\in F}(1-p^{-\beta}).
\]

\begin{lemma}[Valuations and reconstruction]\label{bc:lem:finite}
Let \(\nu\) satisfy \eqref{bc:eq:scaling}. Then
\begin{equation}\label{bc:eq:valuations}
 \nu\{v_p(x)=k_p\text{ for all }p\in F\}
  =\prod_{p\in F}(1-p^{-\beta})p^{-k_p\beta}
   \quad(k_p\in\mathbb N_0).
\end{equation}
In particular, \(\nu(W_F)=c_F(\beta)>0\).
With \(\theta_F=c_F(\beta)^{-1}\nu|_{W_F}\), one has
\begin{equation}\label{bc:eq:finite-reconstruction}
 \int f\,d\nu=c_F(\beta)
   \sum_{n\in\bcN_F}n^{-\beta}\int_{W_F}f(nx)\,d\theta_F(x)
   \quad(f\in C(\bcZhat)).
\end{equation}
The series is absolutely convergent.
\end{lemma}

\begin{proof}
The intersection of \(p^{k_p}\bcZhat\), for \(p\in F\), is
\((\prod_{p\in F}p^{k_p})\bcZhat\). Its measure is
\(\prod_{p\in F}p^{-k_p\beta}\) by scaling.
Finite inclusion-exclusion proves \eqref{bc:eq:valuations}.
For each prime \(p\), the set with zero \(p\)-component has measure
\(\lim_kp^{-k\beta}=0\). Outside this null set for \(p\in F\),
every \(x\) lies in exactly one set \(nW_F\), \(n\in\bcN_F\).
Scaling on each set gives \eqref{bc:eq:finite-reconstruction}.
The total weights sum to one by \eqref{bc:eq:finite-trace}.
\end{proof}

\begin{proposition}[Existence and the invariant measure]\label{bc:prop:radial}
For every \(\beta>0\), there is a unique \(\bcUnits\)-invariant
probability measure \(\nu_\beta\) satisfying \eqref{bc:eq:scaling}.
Under \(\bcZhat=\prod_p\mathbb Z_p\), it is the product of the local measures
\begin{equation}\label{bc:eq:radial}
 \nu_{\beta,p}=(1-p^{-\beta})
       \sum_{k\geq0}p^{-k\beta}(u\mapsto p^ku)_*h_p,
\end{equation}
where \(h_p\) is probability Haar measure on \(\mathbb Z_p^\times\).
\end{proposition}

\begin{proof}
Each local series has total mass one, is invariant under multiplication
by a local unit, and satisfies
\(\nu_{\beta,p}(pB)=p^{-\beta}\nu_{\beta,p}(B)\).
It gives zero mass to zero. Consistent finite products define a positive
norm-one functional on the algebra of continuous cylinder functions.
Uniform extension and Riesz representation produce the product measure.
For an integer \(n\), multiplication by \(n\) is a local unit action
times multiplication by \(p^{v_p(n)}\) at each prime. Only finitely many
of these valuations are nonzero. This proves \eqref{bc:eq:scaling} first
on cylinders and then on all Borel sets by uniqueness of measures.

For uniqueness, fix \(N=\prod_{p\mid N}p^{a_p}\). The action of
\((\mathbb Z/N\mathbb Z)^\times\) on \(\mathbb Z/N\mathbb Z\) has
orbits determined by the truncated valuations
\(\min(v_p(x),a_p)\). The Chinese remainder theorem proves this
description and the surjectivity of \(\bcUnits\) onto the finite unit group.
Invariance makes the measure uniform within each orbit.
Equation~\eqref{bc:eq:valuations} and its tail sums determine every orbit
mass. Therefore every residue class modulo \(N\) has a fixed measure.
These cylinders determine a Borel probability measure on \(\bcZhat\).
\end{proof}

For reference, if a residue \(a\bmod p^j\) has valuation \(k<j\), then
\begin{equation}\label{bc:eq:cylinder}
 \nu_{\beta,p}(a+p^j\mathbb Z_p)
 =\frac{(1-p^{-\beta})p^{-k\beta}}{(p-1)p^{j-k-1}},
 \qquad
 \nu_{\beta,p}(p^j\mathbb Z_p)=p^{-j\beta}.
\end{equation}
These are local marginal probabilities. Products of them give the
probabilities for arbitrary moduli.

\subsection{Uniqueness at and above the critical temperature}

The arithmetic input is a divergence statement for characters of a
finite unit group. The argument only uses Euler products where their
series converge absolutely.

\begin{lemma}[Character divergence]\label{bc:lem:divergence}
Let \(\chi\) be a nontrivial complex character of
\((\mathbb Z/N\mathbb Z)^\times\). Then
\begin{equation}\label{bc:eq:divergence}
 \sum_{p\nmid N}\frac{1-\operatorname{Re}\chi(p)}p=\infty.
\end{equation}
Consequently the same sum with \(p^{-\beta}\) in place of \(p^{-1}\)
diverges whenever \(0<\beta\leq1\).
\end{lemma}

\begin{proof}
Extend \(\chi\) to the integers by zero on nonunits modulo \(N\).
Its sum over a period is zero: multiplication by a unit with nontrivial
character value permutes that sum and multiplies it by that value.
Hence \(A(t)=\sum_{n\leq t}\chi(n)\) satisfies \(|A(t)|\leq N\).
For real \(s>1\), summation by parts gives
\[
 L(s,\chi)=\sum_{n\geq1}\chi(n)n^{-s}
   =s\int_1^\infty A(t)t^{-s-1}\,dt,
 \qquad |L(s,\chi)|\leq N.
\]
The principal character has series
\[
 L_0(s)=\zeta(s)\prod_{p\mid N}(1-p^{-s})\longrightarrow\infty
 \quad(s\downarrow1).
\]
Here \(\zeta(s)\geq\int_1^\infty t^{-s}\,dt=(s-1)^{-1}\),
and the finite product is bounded below by a positive constant.
Thus \(|L(s,\chi)/L_0(s)|\to0\).

Unique factorization and absolute convergence give Euler products for
both series when \(s>1\). Their logarithms satisfy
\begin{equation}\label{bc:eq:log-ratio}
 \log\left|\frac{L(s,\chi)}{L_0(s)}\right|
 =-\sum_{p\nmid N}\frac{1-\operatorname{Re}\chi(p)}{p^s}
  +\sum_{p\nmid N}\sum_{j\geq2}
     \frac{\operatorname{Re}(\chi(p)^j)-1}{jp^{js}}.
\end{equation}
The second sum is bounded in absolute value uniformly for \(s>1\),
because it is at most
\(2\sum_p\sum_{j\geq2}p^{-j}\leq2\sum_{n\geq2}[n(n-1)]^{-1}\).
If \eqref{bc:eq:divergence} failed, the first sum would also be uniformly
bounded. The ratio could then never tend to zero. This contradiction
proves \eqref{bc:eq:divergence}. Comparison of nonnegative terms proves
the assertion for smaller positive \(\beta\).
\end{proof}

\begin{lemma}[Finite-prime contraction]\label{bc:lem:contraction}
Let \(\nu\) satisfy \eqref{bc:eq:scaling}. Suppose that \(f\) factors
through \(\mathbb Z/N\mathbb Z\) and obeys
\(f(vx)=\chi(v)f(x)\) for every unit \(v\) modulo \(N\).
For every finite set \(F\) of primes not dividing \(N\),
\begin{equation}\label{bc:eq:contraction}
 \left|\int f\,d\nu\right|
 \leq\|f\|_\infty
       \prod_{p\in F}\left|\frac{1-p^{-\beta}}
                              {1-\chi(p)p^{-\beta}}\right|
 \leq\|f\|_\infty
       \exp\left(-\sum_{p\in F}(1-\operatorname{Re}\chi(p))p^{-\beta}\right).
\end{equation}
\end{lemma}

\begin{proof}
All \(n\in\bcN_F\) are units modulo \(N\), so
\(f(nx)=\chi(n)f(x)\). Substitute this equality in
\eqref{bc:eq:finite-reconstruction} and sum the absolutely convergent
geometric series. The normalized measure \(\theta_F\) has mass one,
which proves the first inequality. For \(|z|=1\) and \(0<t<1\),
\[
 \log\left|\frac{1-t}{1-zt}\right|
   =\sum_{j\geq1}\frac{\operatorname{Re}(z^j)-1}{j}t^j
   \leq-(1-\operatorname{Re}z)t.
\]
Every summand is nonpositive. Apply this inequality at
\(z=\chi(p)\), \(t=p^{-\beta}\), and exponentiate.
\end{proof}

\begin{theorem}[Uniqueness for \(0<\beta\leq1\)]\label{bc:thm:high}
For every real \(0<\beta\leq1\), the system \((\bcA,\sigma)\) has
exactly one KMS\(_\beta\) state. It is
\(\varphi_\beta=\nu_\beta\circ E\), with \(\nu_\beta\) given by
\eqref{bc:eq:radial}.
\end{theorem}

\begin{proof}
Let \(\nu\) be any scaling measure and fix a modulus \(N\).
The space of complex functions on \(\mathbb Z/N\mathbb Z\) decomposes
into isotypic subspaces for its finite unit group \(G_N\).
Explicitly, its character projections are
\[
 (P_\chi f)(x)=\frac1{|G_N|}\sum_{v\in G_N}
                         \overline{\chi(v)}f(vx),
 \qquad f=\sum_\chi P_\chi f.
\]
Character orthogonality proves the second identity, and
\((P_\chi f)(vx)=\chi(v)(P_\chi f)(x)\).
For nontrivial \(\chi\), Lemmas~\ref{bc:lem:divergence}
and~\ref{bc:lem:contraction} imply \(\nu(P_\chi f)=0\), by exhausting
the primes outside \(N\). Only the trivial character survives.
Thus \(\nu\) is invariant under multiplication by units on every
finite quotient. Locally constant functions are uniformly dense in
\(C(\bcZhat)\), so \(\nu\) is \(\bcUnits\)-invariant.
Proposition~\ref{bc:prop:radial} gives \(\nu=\nu_\beta\).
Theorem~\ref{bc:thm:measures} supplies existence and uniqueness of the state.
\end{proof}

At \(\beta=1\), formula~\eqref{bc:eq:cylinder} assigns mass \(p^{-j}\)
to every residue modulo \(p^j\). Hence \(\nu_1\) is additive Haar
measure on \(\bcZhat\). In particular,
\begin{equation}\label{bc:eq:critical-phase}
 \varphi_1(e(r))=0\quad(r\ne0),\qquad \varphi_1(1)=1.
\end{equation}
The contraction proof establishes the endpoint directly. It requires
neither nonvanishing of \(L(1,\chi)\) nor a factor classification.

\subsection{Every low-temperature state is a Gibbs mixture}

\begin{theorem}[Classification for \(\beta>1\)]\label{bc:thm:low}
Fix \(\beta>1\). For \(\lambda\in\bcProb(\bcUnits)\), define
\begin{equation}\label{bc:eq:gibbs-measure}
 \nu_{\beta,\lambda}
  =\frac1{\zeta(\beta)}\sum_{n\geq1}n^{-\beta}(u\mapsto nu)_*\lambda,
 \qquad \varphi_{\beta,\lambda}=\nu_{\beta,\lambda}\circ E.
\end{equation}
This is an affine homeomorphism from \(\bcProb(\bcUnits)\), with its weak
topology, onto all KMS\(_\beta\) states, with their weak-star topology.
Consequently these states form a simplex. Its extreme points are
\(\varphi_{\beta,u}=\varphi_{\beta,\delta_u}\), \(u\in\bcUnits\), and
\begin{equation}\label{bc:eq:gibbs-state}
 \varphi_{\beta,u}(a)
   =\zeta(\beta)^{-1}\bcTr\big(\pi_u(a)e^{-\beta H}\big).
\end{equation}
\end{theorem}

\begin{proof}
Let \(\nu\) be any scaling measure. Since
\(\bcUnits=\bigcap_p(\bcZhat\setminus p\bcZhat)\), continuity from above and
Lemma~\ref{bc:lem:finite} give
\[
 \nu(\bcUnits)=\prod_p(1-p^{-\beta})=\zeta(\beta)^{-1}.
\]
The Euler product equality follows by expanding each finite product and
then exhausting the positive integers. Absolute convergence holds for
\(\beta>1\). The sets \(n\bcUnits\), \(n\geq1\), are pairwise disjoint:
their valuation vectors determine \(n\). By scaling,
\[
 \sum_{n\geq1}\nu(n\bcUnits)
 =\nu(\bcUnits)\sum_{n\geq1}n^{-\beta}=1.
\]
Their union therefore has full measure. Define
\(\lambda=\zeta(\beta)\nu|_{\bcUnits}\), a probability measure on the
compact set \(\bcUnits\). Applying scaling to every Borel subset of
\(\bcUnits\) proves \eqref{bc:eq:gibbs-measure}. This also proves uniqueness
of the parameter measure.

Conversely, the series in \eqref{bc:eq:gibbs-measure} defines a probability
measure. The point \(nu\) lies in \(m\bcZhat\) exactly when \(m\mid n\),
because \(u\) is a unit at every prime. Thus its mass on \(mB\) is
\[
 \frac1{\zeta(\beta)}\sum_{k\geq1}(mk)^{-\beta}
          \lambda\{u:ku\in B\}
 =m^{-\beta}\nu_{\beta,\lambda}(B).
\]
Theorem~\ref{bc:thm:measures} now gives all the claimed states.

For each \(f\in C(\bcZhat)\), the series
\(\zeta(\beta)^{-1}\sum_n n^{-\beta}f(nu)\) converges uniformly in
\(u\). It is continuous, which proves continuity of the affine map
on parameter measures. A continuous bijection from a compact space to
a Hausdorff space is a homeomorphism. The extreme points of
\(\bcProb(\bcUnits)\) are precisely its point masses: a measure that is not
a point mass has a Borel set of mass strictly between zero and one,
which gives a nontrivial convex decomposition. Every measure has the
unique decomposition \(\lambda=\int\delta_u\,d\lambda(u)\).
This proves the simplex assertion without invoking a general KMS
decomposition theorem.

Finally, on \(f\in D\) the normalized trace in
\eqref{bc:eq:gibbs-state} has value
\(\zeta(\beta)^{-1}\sum_nn^{-\beta}f(nu)\).
On \(\mu_mf\mu_n^*\) with \(m\ne n\), every diagonal matrix
coefficient is zero: a nonzero coefficient sends index \(k\) to
\(mk/n\), which differs from \(k\). For \(m=n\), the element is
\(\alpha_n(f)\in D\). The trace state and \(\varphi_{\beta,u}\)
therefore agree on a dense algebra and hence everywhere.
\end{proof}

\subsection{The cyclotomic carrier and the symmetry action}

Let \(\mathbb Q^{\mathrm{cycl}}\subset\mathbb C\) be the union of all
fields \(\mathbb Q(\zeta_N)\), with \(\zeta_N=\exp(2\pi i/N)\).
The group \(G=\bcGal(\mathbb Q^{\mathrm{cycl}}/\mathbb Q)\) carries
the inverse-limit topology of its finite cyclotomic Galois groups.

\begin{lemma}[Cyclotomic identification]\label{bc:lem:galois}
The map \(u\mapsto\gamma_u\), specified by
\begin{equation}\label{bc:eq:galois}
 \gamma_u(\zeta_N)=\zeta_N^{u_N},
\end{equation}
is a topological group isomorphism \(\bcUnits\longrightarrow G\).
\end{lemma}

\begin{proof}
The cyclotomic polynomial \(\Phi_N\) is monic with integer coefficients,
as follows inductively from \(X^N-1=\prod_{d\mid N}\Phi_d(X)\).
We recall the irreducibility argument to check the full parameter space.
Let \(f\in\mathbb Z[X]\) be the monic irreducible polynomial of a
primitive root \(\zeta_N\). It divides \(\Phi_N\), by Gauss's lemma.
If a prime \(p\nmid N\) sent this root to a root of a different
irreducible factor \(g\) of \(\Phi_N\), then
\(f(X)\mid g(X^p)\). Reduction modulo \(p\) gives
\(\overline f\mid\overline g^p\), so \(\overline f\) and
\(\overline g\) have a common irreducible factor.
It would occur twice in \(\overline\Phi_N\) and thus in \(X^N-1\)
over \(\mathbb F_p\). This is impossible because the derivative
\(NX^{N-1}\) is coprime to \(X^N-1\).

The same argument applies to each root of \(f\). Consequently every
map \(\zeta_N\mapsto\zeta_N^p\), \(p\nmid N\), preserves its roots.
Every positive integer coprime to \(N\)
is a product of such primes. Thus all primitive roots are roots of
\(f\), and \(\Phi_N=f\). It follows that each unit modulo \(N\)
defines a field automorphism by \eqref{bc:eq:galois}.
Compatible residues define compatible automorphisms, hence an
automorphism of their union. Conversely, every automorphism preserves
root orders and determines these compatible units. This proves the
bijection and the group law. The finite-quotient definitions of the
two topologies prove continuity in both directions.
\end{proof}

Under this identification, \eqref{bc:eq:representation} becomes
\[
 \pi_\gamma(\mu_m)\varepsilon_n=\varepsilon_{mn},\qquad
 \pi_\gamma(e(r))\varepsilon_n
      =\gamma(e^{2\pi i nr})\varepsilon_n.
\]
Theorem~\ref{bc:thm:low} therefore has exactly the cyclotomic Galois
carrier in the Bost--Connes classification
\cite{BostConnes1995}*{Propositions~23--24 and Theorem~25}.

\begin{theorem}[Bost--Connes equilibrium classification]
\label{thm:bc}
The full system \((\bcA,\sigma)\) has a unique KMS\(_\beta\) state for
\(0<\beta\leq1\). For \(\beta>1\), its KMS simplex is affinely
homeomorphic to \(\bcProb(\bcUnits)\), and its extreme states are the
normalized Gibbs states \eqref{bc:eq:gibbs-state}.
Via \eqref{bc:eq:galois}, their parameters are exactly
\(\operatorname{Gal}(\mathbb Q^{\mathrm{cycl}}/\mathbb Q)\).
\end{theorem}

\begin{proof}
Theorems~\ref{bc:thm:measures}, \ref{bc:thm:high}, and~\ref{bc:thm:low}
construct and classify the states. Lemma~\ref{bc:lem:galois} identifies
the parameter group. In particular, no normalization of a divergent
occupation trace is used at or below \(\beta=1\).
\end{proof}

\begin{proposition}[Symmetry and the Toeplitz restriction]\label{bc:prop:symmetry}
For \(v\in\bcUnits\), the formulas
\[
 \theta_v(\mu_n)=\mu_n,\qquad \theta_v(e(r))=e(vr)
\]
define an action of \(\bcUnits\) on \(\bcA\) by automorphisms that commute
with \(\sigma\). Here \(vr\) uses the residue of \(v\) modulo the
order of \(r\). Its action on the low-temperature extreme states is
\begin{equation}\label{bc:eq:symmetry}
 \varphi_{\beta,u}\circ\theta_v=\varphi_{\beta,vu},
 \qquad \beta>1.
\end{equation}
It is free and transitive on those states. For \(0<\beta\leq1\),
the unique KMS state is fixed by every \(\theta_v\).

For every \(\beta>0\), all KMS\(_\beta\) states have the same
restriction to \(\mathcal T=C^*(\mu_n:n\in\bcN)\). On its dense
monomials this restriction is
\begin{equation}\label{bc:eq:toeplitz}
 \varphi(\mu_m\mu_n^*)=\delta_{m,n}n^{-\beta}.
\end{equation}
On any finite set of primes it is the product of the geometric
Toeplitz states with weights \((1-p^{-\beta})p^{-k\beta}\).
\end{proposition}

\begin{proof}
Multiplication by a profinite unit permutes \(\mathbb Q/\mathbb Z\)
and each fiber of multiplication by \(n\). Hence \(\theta_v\)
preserves all the defining relations. Its inverse is
\(\theta_{v^{-1}}\), and it commutes with the dynamics on generators.
The identity \(\pi_u\circ\theta_v=\pi_{vu}\) gives
\eqref{bc:eq:symmetry}. The parameterization in
Theorem~\ref{bc:thm:low} proves freeness and transitivity.
At \(0<\beta\leq1\), invariance also follows from
Theorem~\ref{bc:thm:high} or from the invariant measure.

The KMS identity gives zero at unequal indices and \(n^{-\beta}\)
at equal indices. The monomials \(\mu_m\mu_n^*\) span a dense star
algebra in \(\mathcal T\), by the product calculation in
Lemma~\ref{bc:lem:expectation}. Thus \eqref{bc:eq:toeplitz} determines the
restriction. For finitely many primes, unique factorization identifies
the diagonal powers and gives the product of their local values.
Those local values are the geometric-state expectations, since
\(\sum_{k\geq j}(1-p^{-\beta})p^{-k\beta}=p^{-j\beta}\).
\end{proof}

There is exactly one symmetry-invariant KMS state for each \(\beta>0\).
For \(\beta>1\), it is the mixture in \eqref{bc:eq:gibbs-measure} with
probability Haar measure on \(\bcUnits\). It agrees with
\(\nu_\beta\circ E\) by Proposition~\ref{bc:prop:radial}.
For a nonzero \(r=a/b\) in lowest terms, its phase value is
\begin{equation}\label{bc:eq:phase-value}
 \varphi_\beta(e(r))
   =\prod_{p^j\Vert b}
         p^{-j\beta}\frac{1-p^{\beta-1}}{1-p^{-1}}.
\end{equation}
Indeed, the average of a primitive additive character of conductor
\(p^j\) on \(p^k\mathbb Z_p^\times\) is zero for \(k<j-1\),
is \(-1/(p-1)\) for \(k=j-1\), and is one for \(k\geq j\).
The first assertion follows by subtracting the complete additive sum
on multiples of \(p\) from the complete sum modulo \(p^j\).
Inserting these three values in \eqref{bc:eq:radial} gives the local factor.
For example, \(\varphi_\beta(e(1/2))=2^{1-\beta}-1\).

\subsection{Critical and zero-temperature limits}

\begin{proposition}[Limits of the equilibrium states]\label{bc:prop:limits}
For every \(a\in\bcA\),
\begin{equation}\label{bc:eq:critical-limit}
 \sup_{\lambda\in\bcProb(\bcUnits)}
       |\varphi_{\beta,\lambda}(a)-\varphi_1(a)|\longrightarrow0
       \quad(\beta\downarrow1).
\end{equation}
At zero temperature, define
\[
 \varphi_{\infty,u}(a)
  =\langle\pi_u(a)\varepsilon_1,\varepsilon_1\rangle,
 \qquad
 \varphi_{\infty,\lambda}=\int_{\bcUnits}\varphi_{\infty,u}\,d\lambda(u).
\]
These are states, and
\begin{equation}\label{bc:eq:zero-limit}
 \|\varphi_{\beta,\lambda}-\varphi_{\infty,\lambda}\|
       \leq2\bigl(1-\zeta(\beta)^{-1}\bigr)\longrightarrow0
       \quad(\beta\to\infty),
\end{equation}
uniformly in \(\lambda\). The weak-star limits of KMS states as
\(\beta\to\infty\) are exactly these mixtures. Their extreme points
are indexed by \(\bcUnits\), or equivalently by \(G\), and
\(\varphi_{\infty,u}(e(r))=\bcphase_r(u)\).
\end{proposition}

\begin{proof}
Suppose \eqref{bc:eq:critical-limit} fails. There are \(\beta_j\downarrow1\)
and parameter measures \(\lambda_j\) whose state values stay a fixed
positive distance from \(\varphi_1(a)\). The state space is weak-star
compact, so a subnet converges to a state \(\psi\). For
\(b,c\in\mathcal B\), the KMS identity passes to the limit because
\(\sigma_{i\beta_j}(b)\to\sigma_i(b)\) in norm. Therefore \(\psi\)
is KMS\(_1\) by Lemma~\ref{bc:lem:analytic}, and
Theorem~\ref{bc:thm:high} gives \(\psi=\varphi_1\). This contradicts
the choice of the subnet.

For the other limit, the Gibbs formula writes each state as a convex
sum of vector states at \(\varepsilon_n\). The weight at
\(\varepsilon_1\) is \(\zeta(\beta)^{-1}\). The remaining weights
have total mass \(1-\zeta(\beta)^{-1}\). Two states are at norm distance
at most two, which proves \eqref{bc:eq:zero-limit} and its integrated form.
The zeta series tends to one by dominated convergence, using any fixed
exponent greater than one as a summable majorant.

The map \(u\mapsto\varphi_{\infty,u}(a)\) is continuous first on
monomials and then on all \(a\) by uniform approximation. Thus mixtures
depend continuously on \(\lambda\). Compactness of \(\bcProb(\bcUnits)\)
and \eqref{bc:eq:zero-limit} identify every weak-star limit with such a
mixture. Conversely, keeping \(\lambda\) fixed produces every mixture.
The phase algebra separates measures on \(\bcUnits\), since its restricted
characters are uniformly dense in \(C(\bcUnits)\). Hence the mixture map
is injective and its extreme points are precisely the point measures.
\end{proof}

The zero-temperature assertion concerns limits of finite-temperature
KMS states; it does not classify all ground states of the dynamics.
The normalization in \eqref{bc:eq:gibbs-state} always uses the
trace-class half-plane \(\operatorname{Re}s>1\). The state at the
critical point is instead constructed by the scaling measure, and
\eqref{bc:eq:critical-limit} specifies its relation to the Gibbs family.


\subsection{Vacua and augmentations}

The occupation vacuum is the vector \(\varepsilon_1\), of energy zero.
Its vector state \(\varphi_{\infty,u}\) is a zero-temperature limit,
whereas a KMS\(_\beta\) state with finite \(\beta>0\) obeys
\(\varphi(\mu_n\mu_n^*)=n^{-\beta}\).
An algebra augmentation means a unital multiplicative map to the
coefficient field. Here the formulas
\[
 \epsilon(e(r))=1,\qquad \epsilon(\mu_n)=1
\]
preserve every defining relation and give a star augmentation
\(\epsilon:\bcA\to\mathbb C\). It has
\(\epsilon(\mu_n\mu_n^*)=1\), so it is not KMS\(_\beta\) for
\(\beta>0\). The vacuum state is not multiplicative either: for
\(n>1\), its value on \(\mu_n\) is zero, while its value on
\(\mu_n^*\mu_n\) is one. Thus the equal scalar Euler factor in the
local integral and occupation trace specifies none of these three
additional choices. A chiral differential or arithmetic cohomological
comparison must be constructed on further data.
